\documentclass[a4paper,11pt]{article}

\usepackage{lscape}
\usepackage{geometry}
\usepackage[utf8]{inputenc}
\usepackage[francais]{babel}
\usepackage[T1]{fontenc}
\usepackage{gantt}
\usepackage{tikz}
\usepackage{tikz-uml}
\usepackage{relsize}
\usepackage{color}
\usepackage[hidelinks]{hyperref}
\definecolor{dkgreen}{rgb}{0,0.6,0}
\definecolor{gray}{rgb}{0.5,0.5,0.5}
\definecolor{mauve}{rgb}{0.58,0,0.82}

\usepackage{listings}
\usepackage{float}
%\usepackage{kpfonts}

\usepackage{graphicx}
%\usepackage{rotating}

\lstset{
  language=C++,
  basicstyle=\footnotesize,
  backgroundcolor=\color{white},
  keywordstyle=\color{red},
  commentstyle=\color{dkgreen},
  stringstyle=\color{mauve},
  numberstyle=\color{red},
  morekeywords={string},
  frame=BL,
  aboveskip=1em,
  belowskip=2em,
}
\lstset{
  literate={ù}{{\`u}}1
  {é}{{\'e}}1
  {è}{{\'e}}1
  {à}{{\`a}}1
}


\lstdefinelanguage{tikzuml}{language=[LaTeX]TeX, classoffset=0, morekeywords={umlbasiccomponent, umlprovidedinterface, umlrequiredinterface, umldelegateconnector, umlassemblyconnector, umlVHVassemblyconnector, umlHVHassemblyconnector, umlnote, umlusecase, umlactor, umlinherit, umlassoc, umlVHextend, umlinclude, umlstateinitial, umlbasicstate, umltrans, umlstatefinal, umlVHtrans, umlHVtrans, umldatabase, umlmulti, umlobject, umlfpart, umlcreatecall, umlclass, umlvirt, umlunicompo, umlimport, umlaggreg}, classoffset=1, morekeywords={umlcomponent, umlsystem, umlstate, umlseqdiag, umlcall, umlcallself, umlfragment, umlpackage}, classoffset=0,  sensitive=true, morecomment=[l]{\%}}

\geometry{margin=2.5cm}
\geometry{headheight=15pt}

\usepackage{fancyhdr}
\usepackage{fancyvrb}
\usepackage{float}
\usepackage[footnote,smaller]{acronym}

\pagestyle{fancy}
\rhead{RE203 - Projet de Réseau}

% \acrodef{LABRI}{Laboratoire Bordelais de Recherche en Informatique}

\begin{document}

\begin{titlepage}
  \begin{center}

    \textsc{\LARGE RE203 - Projet de Réseau}\\[2cm]
    \textsc{\large Rapport Final}\\[3cm]
    Maxime \textsc{Bellier} \ \ \ Jean-Michaël \textsc{Celerier}\\
    Julien \textsc{Chaumont} \ \ \ Bazire \textsc{Houssin} \ \ \ Sylvain \textsc{Vaglica}\\[1cm]
    \textsc{Groupe 3}\\[1.5cm]
    \textsc{\large 23/05/2013 }\\[1.5cm] %TODO
    \includegraphics[width=8cm]{logo.png}

  \end{center}
  \vspace{3cm}

\end{titlepage}

\clearpage

\section*{Introduction}

Ce rapport final pour le projet de réseau RE203 fait suite au rapport intermédiaire délivré le 23 avril dernier. Il sera ici question de l'implémentation mise en \oe uvre par notre groupe, l'explication du sujet ayant été développée précédemment.

\section{Le bloc commun}

Étant donné que le développement des deux parties du projet (contrôleur et routeur) a été réalisé respectivement dans les langages C et C++, il a été décidé d'établir une sorte bibliothèque commune sur laquelle se baser. Cette partie comprend:
\begin{itemize}
 \item une invite de commande
 \item un système de gestion des fichiers de configuration
 \item des fonctions de traitement des messages reçus et à envoyer
 \item une API pour la gestion des sockets
 \item quelques fonctions utilitaires
\end{itemize}

Chacun de ces blocs est indépendant des autres, ce qui assure une intégration facilitée dans le contrôleur comme dans le routeur.

\subsection{Le prompt}

\underline{Fichiers :} \texttt{common/prompt.c}, \texttt{common/prompt.h}\\

L'invite de commande, ou \textit{prompt} en Anglais, permet la communication entre le contrôleur ou le routeur, et l'utilisateur. Il suffit pour cela de regarder constamment sur l'entrée standard et de traiter les commandes saisies par l'utilisateur. Afin de ne pas bloquer le programme sur l'écoute du fichier \texttt{stdin}, il faut pouvoir exécuter cette tâche en parallèle, d'où l'utilisation de la bibliothèque \texttt{pthread} ici.

Les fonctions définies dans les fichiers sus-mentionnés assurent la création de ce thread et l'écoute sur l'entrée standard. Le paramètre à fournir à la fonction d'initialisation est un pointeur de fonction. Lorsqu'une information arrive, elle est mise en mémoire et convertie en un \texttt{Message}\footnote{Se référer à la partie \ref{message}}, et c'est à la fonction dont le pointeur est passé en paramètre de gérer les actions à effectuer suite à la réception de ce message.

De cette manière, contrôleur et routeur ont seulement à gérer la fonction de traitement des messages utilisateur (chaque programme devant comprendre des commandes qui lui sont spécifiques), mais la tâche de récupération des instructions de l'utilisateur est quant à elle unifiée.

\subsection{Traitement des fichiers de configuration}

\underline{Fichiers :} \texttt{common/config.c}, \texttt{common/config.h}\\

Le sujet demandait à ce que routeur comme contrôleur puissent être paramétrés depuis un fichier de configuration. Ce module permet la lecture de ces fichiers. La bibliothèque C \texttt{T-Rex}\footnote{Publiée sous licence libre \textit{zlib}, consultable à l'adresse \url{http://sourceforge.net/projects/tiny-rex/}}, qui permet la manipulation d'expressions régulières, a été utilisée dans ce cadre afin de faciliter la reconnaissance des différentes lignes possibles dans les fichiers de configuration.

Dans le code du routeur (respectivement du contrôleur), il suffit d'appeler la fonction \texttt{config\_\_readRouter} (respectivement \texttt{config\_\_readController}). En interne, ces fonctions ouvrent le fichier de configuration indiqué, le lisent ligne par ligne, identifient les mots-clés à l'aide des expressions régulières, et remplissent une variable de type \texttt{Configuration} avec les informations contenues dans le fichier. Cette structure de données est déclarée de la manière suivante:
\begin{lstlisting}
enum SoftwareType { ROUTER, CONTROLLER };

typedef struct Configuration
{
	enum SoftwareType type;
	char controllerAddress[20];
	int routerPort;
	int controllerPort;
	int pollTimeoutValue;
	int controllerUpdateInterval;
	int routerUpdateInterval;
	int defaultTTLValue;
	int defaultPingPacketCount;
	int defaultPacketTimeoutValue;
	int defaultDVTimeoutValue;
} Configuration;
\end{lstlisting}

A la fin de la lecture, il ne reste plus qu'à retourner au programme appelant un pointeur vers la structure remplie, qui n'aura alors qu'à consulter la valeur de ses différents attributs.

Notez que la structure \texttt{Configuration} contient des attributs concernant aussi bien le contrôleur que le routeur. C'est au programme de veiller à n'utiliser que les attributs le concernant, les autres attributs n'ayant aucune garantie d'avoir été initialisés à une valeur cohérente (précisons cependant qu'il n'y a aucune raison pour que le contrôleur ait besoin des informations propres aux routeurs, et inversement).

\section{Le contrôleur}


\section{Les routeurs}


\end{document}
